\documentclass{jsarticle}
\usepackage[dvipdfmx,final]{graphicx}
\usepackage{amsmath}
\usepackage{amssymb}
\begin{document}
\title{}
\author{}
\date{}
\maketitle
\thispagestyle{empty}

    \newcommand{\variint}{%
	\begin{picture}(15,30)
		\put(0,0){
			$ \iint $}
		\put(0,5){\line(15,0){15}}
	\end{picture} %
}

\newcommand{\variiint}{%
	\begin{picture}(15,15)
		\put(0,0){
			$ \iiint $}
		\put(0,5){\line(15,0){20}}
	\end{picture} %
}



 \newcommand{\alearletter}{%
	 \begin{picture}(10,20)
		 \put(0,0){
			 $ \square $}
		 \put(0,0){\line(1,1){10}}
	 \end{picture} %
	 }

\newcommand{\letter}{%
	\begin{picture}(5,5)
		\put(0,0){
			$ \square $}
		\put(3,3){$\vee$}
	\end{picture} %
}
 \newcommand{\secproduct}{%
	 \begin{picture}(10,20)
		 \put(0,0){
			 $ \nabla $}
		 \put(5,0){+}
	 \end{picture} %
	 }
 \newcommand{\average}{%
	 \begin{picture}(5,5)
		 \put(0,0){
			 $ \square $}
		 \put(3,3){$\slash$}
	 \end{picture} %
	 }
\newcommand{\averagesqrt}{%
	 \begin{picture}(5,5)
		 \put(0,0){
			 $ \square $}
		 \put(3,3){$\setminus $}
	 \end{picture} %
	 }




新年あけまして、おめでとうございます。
今年もよろしくお願いいたします。
私のこと、捨てないでください。
障害者は、事情で仕方なかったのです。
金星の酉の手相ですから、結果オーライです。
いつか、もとにもどります。
だから、見捨てないでください。
家具屋の社長さん、同志社大学の先生に、山口家背負って立つでしょう
と言われているのを覚えています。
1995年は本当に、申し訳ございませんでした。
以下は、アカシックレコードからの、ロールシャッハテストの能力で
見つけている数式たちです。
私、切り離されても大丈夫です。
陛下が、いつか上皇様と同じことをするらしく、愛子様、
皇族で勝利です。
苫米地博士の夢は、この世から、病気を無くすことと、博士の本で知っています。
以下の式たちは、数式のサラダではないです。
ロールシャッハテストと、数式の製造で出来ています。



  $$\int_M w dw_x = \int_{D\chi}|_u dx_m, 
   \mathcal{O}|_M = \mathcal{O}_{D\chi},, 
   w = (1 + x)^{1 \over 2}, 
   (1 + i)^n = \int z dz_m $$
  $$ \boxtimes \to x \boxtimes y \to (x \cdot y)^{\times}, 
   \lambda_{[xy]_m} = [\lambda_x \otimes \lambda_y]^{\times z} $$

  $$ \letter = (\nabla \oplus \square)^{\nabla L}, 
   = (\nabla \otimes \square)|^{L}_{D\chi}, 
   \ll \secproduct|_x, \secproduct|_y \gg $$
  $$ = \boxplus^{\nabla L} $$
  $$ H\Psi = i \hbar \psi, 
   = [\vartriangleright|^{\times L_n}, 
  \blacktriangleright|^{\times L_m}], 
   V_{D\chi}\int dx_m = S \otimes S, 
   = V\int dx_m = S^2 \oplus S^2 
   \int R^{\nabla r}dr_x = l^{R|_{\nabla}}, 
   = \beta(p,q), 
   x \boxtimes y = [\average x\cdot y, \averagesqrt x \cdot y] $$ 

 
	 $$ \square \Psi = {}^{t}\scalebox{1}[2]{\variiint}
	 \mathrm{cohom}D_{\chi}(\chi,x)[Im], 
  {}^{t}\scalebox{1}[2]{\variiint}\mathrm{cohom D_{\chi}[I_m]} $$
$$ = ||ds^2|| $$
$$ \mathcal{\psi}_{\mu\nu} = {\partial \over {\partial \Psi}} {}^{t}
 \scalebox{1}[2]{\variiint} \mathcal{\psi}(x,y,z)dm, 
  \int (T^{\mu\nu})^{'}dx_m = \int (R + {1 \over 2}\Lambda g_{ij})dx_m $$
 $$ = \int e^{x\log x}\cdot \mathrm{div}(\mathrm{rot} E)dx, 
  = e^{-x \log x} $$
 $$ \square_{\mu\nu} = \bigoplus^{\infty}_{\mu\nu} \mathcal{\psi}_{\mu\nu}
 (x,y,z)d\Psi, 
  \int e^{-t}x^{1 - t}dx = \bigoplus a^{tx}x^{t}[I_m] \to a^{tx}x^{t -1} $$
 $$ = \bigoplus {(i \hbar^{\nabla})}^{\oplus L} $$
 $$ {d \over {d l}}L(x,y) = \int [D^2 \psi \otimes h\nu]d \tau, 
  \square_{\mu\nu} = \bigoplus \mathcal{\psi}_{\mu\nu}(x,y,z)d\Psi $$
 $$ \tau(k) = \mathcal{N}\nu(ij) \nabla_{ij}\sum a_k f^{\mu\nu}, 
  = \square \mathcal{\psi}_{\mu\nu}(x_m) $$
 $$ ||ds^2|| = \cap_{k=0}\mathcal{\psi}_{\mu\nu}(I_m) 
  D\int g_{ij}|^{\oplus L_{ij}}_{\nu(\tau)} = \nabla \nabla_{ij}\int 
 \nabla f(\bar{x})\cdot x d\eta $$
 
 \end{document}
